% ARCHIVED DRAFT: independently audited by root and incorporated on 2026-11-19 UTC.
% The assembled paper is authoritative; this file retains drafting provenance.
\begin{saresult}{B.074}{Proved under stated hypotheses}{Regular sequences at finite-residue points without Noetherianity}
Let $R$ be a nonzero commutative ring such that, for every maximal ideal
$\mathfrak m$ with finite residue field, the maximal ideal of
$R_{\mathfrak m}$ contains a regular sequence of length two.
For every unital extension $R\subset S$, the polynomial co-operations
exist uniquely exactly when
\[
 (R:_S D_d(R))=R\quad\text{for every }d\ge1,
 \quad\text{equivalently}\quad
 \operatorname{Ann}_{S/R}(D_d(R))=0\quad\text{for every }d. \tag{B74.1}
\]
This is also equivalent to tensor injectivity, to $C=R$, and to
$\Omega(R,S)=R[X]$. All positive cases have all tensor isomorphisms.
There is no Noetherian, reducedness or coefficient-torsion hypothesis.
The proof does not presume localization of $\Omega$.
\end{saresult}
\begin{proof}
Use the actual coefficient ring $C$ of B.064. If the second tensor map
is injective, $R\to C$ is an injective epimorphism. Localize that map
at a maximal ideal to obtain an injective epimorphism $A\subset B$.
The exact determinant-ideal localization (B69.3), which holds for
every ring, implies that each $b\in B$ is multiplied into $A$ by some
$D_d(A)$. If the residue field is infinite, that determinant ideal is
$A$, so $B=A$.

Suppose the residue field is finite, and let $a,t$ be a regular sequence
in its maximal ideal. The case $d=1$ is immediate. For $d\ge2$ the
explicit polynomial $P_d$ of B.064 gives
$a^{e_d},t^{e_d}\in D_d(A)$, since its values are those powers times
units, as in (B69.4). Write $e=e_d$. The pair $a^e,t^e$ is still a
regular sequence. For clarity, regularity of $t$ modulo $a^j$ follows
by induction: if $tx\in a^jA$, reduction modulo $a$ gives $x=ay$;
cancelling the regular $a$ reduces the assertion to exponent $j-1$.
Iterating regularity of $t$ gives regularity of $t^e$ modulo $a^e$.

Let $W$ denote the regular elements of $A$. The inclusion $A\subset B$
induces an inclusion of the total quotient ring $Q(A)$ in $W^{-1}B$.
For $b\in B$ with $D_d(A)b\subset A$, put $x=a^eb$ and $y=t^eb$,
both in $A$. The equality $t^ex=a^ey$ and regularity modulo $a^e$ show
that $x=a^ez$ for some $z\in A$. Thus the image of $b$ in $W^{-1}B$
is exactly $z$. In particular every rationalized element of $B$ lies
in the embedded copy of $A$.

Restriction of rationalization therefore gives a ring retraction
$\rho:B\to A$ of the inclusion. The two endomorphisms of $B$ consisting
of its identity and of inclusion composed with $\rho$ agree on $A$.
The epimorphism property forces them to agree on $B$, hence $B=A$.
This holds at every maximal ideal, so $C/R=0$ by the arbitrary-module
localization test. B.064 gives precisely the equivalences in (B74.1)
and the positive ordinary-polynomial tensor isomorphisms.
\end{proof}
